\documentclass[twocolumn]{iopjournal}

\usepackage{amsmath,amssymb,amsfonts}
\usepackage{array}

\fancyhfoffset[L]{0mm}
\setlength{\marginparwidth}{0pt}
\setlength{\marginparsep}{0pt}

% ==========================================
% ZENODO PREPRINT HEADERS (ACTIVE)
% ==========================================
\renewcommand{\articletype}[1]{%
  {\vspace*{-8mm}\noindent \Large \sf Preprint}\par
  \vspace*{6mm}\noindent{\scriptsize \sf{\bfseries \MakeUppercase{#1}}}}

\fancyhead[L]{}
\fancyhead[R]{}

\newcommand{\sgn}{\mathrm{sgn}}
\newcommand{\std}{\mathrm{std}}
\newcommand{\Ind}{\mathrm{Ind}}

\hypersetup{pdftitle={The Torsion Sector of the Tetrad Connection as the Regular Representation of S4}, pdfauthor={Stephen Nelson}, pdfsubject={Gravity Field Dynamics (GFD)}, pdfkeywords={torsion, affine connection, representation theory, symmetric group, Einstein-Cartan, proof}}

\begin{document}

\articletype{Supplementary Material}

\title{The Torsion Sector of the Tetrad Connection\\
as the Regular Representation of $S_4$}

\author{Stephen Nelson}

\affil{Field Dynamics Research}

\ead{stephen@fielddynamics.org}

\begin{abstract}
The symmetric group $S_4$ permutes the four legs of a tetrad, and
with them the $64$ components of the affine connection. We show that
the $24$ torsion components form the regular representation of
$S_4$, with an explicit generator: a single torsion tensor whose $24$
images under the permutations of the legs are a basis. The proof
splits the torsion components into two orbits of $12$, according to
whether the upper index repeats a lower one. The sign representation occurs once in the torsion
sector and not at all in the $40$-component Christoffel sector, and
it is spanned by the totally antisymmetric torsion
$\sum_d \varepsilon_{abcd}$. With indices lowered by $|\eta| = I$,
the vector, axial and tensor parts of the torsion lie in the two
halves as $12 = 4 + 8$ and $12 = 4 + 8$, and the two tensor halves
are isomorphic; lowered by $\eta$ itself, the axial part is not
$S_4$-invariant. The regular representation fits the torsion sector
only for $K = 2$ and $K = 4$ legs.
\end{abstract}

\keywords{torsion, affine connection, representation theory, symmetric group, Einstein--Cartan, proof}

% ----------------------------------------------------------------------

\section*{Setting}

Let $V = \mathbb{R}^4$ have basis $f_0, \dots, f_3$, one vector per
leg, and let $S_4$ act by $\sigma f_a = f_{\sigma(a)}$. Read by
magnitude, the legs are equivalent: $|\eta_{ab}| = \delta_{ab}$ is
invariant under every permutation~\cite{Nelson2026frame}. The dual
basis is permuted in the same way, so $V^* \cong V$ as
$S_4$-modules with no metric. The affine connection
$\Gamma^a{}_{bc}$ lies in $V \otimes V^* \otimes V^* \cong V \otimes V
\otimes V$, which splits into
the Christoffel sector $C = V \otimes \mathrm{Sym}^2 V$
($40$ components) and the torsion sector
$T = V \otimes \Lambda^2 V$ ($24$ components), both
$S_4$-submodules~\cite{Hehl1976}. The torsion tensor takes values in
$T$ in any basis; in a coordinate basis, $C$ and $T$ are the symmetric
and antisymmetric parts of $\Gamma$. $T$ has basis
$t_{a;bc} = f_a \otimes f_b \wedge f_c$, $b < c$, and
$\sigma\, t_{a;bc} = t_{\sigma(a);\sigma(b)\sigma(c)}$, with
$t_{a;cb} = -t_{a;bc}$.

We write $\mathbf{1}$, $\sgn$, $[2,2]$, $\std$ and
$\sgn \otimes \std$ for the irreducible representations of $S_4$, of
dimensions $1, 1, 2, 3, 3$~\cite{FultonHarris}, and $H = \langle
(23) \rangle$. Theorems 1 and 2 and the Proposition use only the
labels of the legs; Theorem 3 lowers an index with $|\eta| = I$.

% ----------------------------------------------------------------------

\section*{Two Orbits}

\paragraph{Lemma.} \textit{Up to sign, the basis vectors of $T$ form
two $S_4$-orbits of $12$: the incident vectors, with
$a \in \{b,c\}$, and the non-incident vectors, with
$a \notin \{b,c\}$. The stabiliser of the line of $t_{0;01}$, and of
the line of $t_{0;23}$, is $H$; it acts by $+1$ on $t_{0;01}$ and by
$-1$ on $t_{0;23}$.}

\paragraph{Proof.} An incident vector is determined by the ordered pair
$(a, b)$ of distinct labels with $\{b, c\} = \{a, b\}$, and a
non-incident vector by the ordered pair $(a, d)$, where $d$ is the
label outside $\{a, b, c\}$. $S_4$ is transitive on the $12$ ordered
pairs of distinct labels, so each family is one orbit. A permutation
fixing the line of $t_{0;01}$ fixes $0$ and the set $\{0,1\}$, hence
$1$, and is $e$ or $(23)$; $(23)\,t_{0;01} = t_{0;01}$. One fixing
the line of $t_{0;23}$ fixes $0$ and the set $\{2,3\}$, hence $1$,
and is $e$ or $(23)$; $(23)\,t_{0;23} = t_{0;32} = -t_{0;23}$.
\hfill $\square$

% ----------------------------------------------------------------------

\section*{The Regular Representation}

\paragraph{Theorem 1.} \textit{Let
$\xi = t_{0;01} + t_{0;23}$. The $24$ vectors $\sigma\xi$,
$\sigma \in S_4$, are a basis of $T$, and}
\begin{equation}
  \Phi : \mathbb{R}[S_4] \to T, \qquad \sigma \mapsto \sigma\xi,
  \label{eq:Phi}
\end{equation}
\textit{is an isomorphism of $S_4$-modules: the torsion sector is
the regular representation.}

\paragraph{Proof.} With $\tau = (23)$, the Lemma gives
$\tau\xi = t_{0;01} - t_{0;23}$. Hence, for every $\sigma$,
$\sigma\xi + \sigma\tau\xi = 2\,\sigma t_{0;01}$ and
$\sigma\xi - \sigma\tau\xi = 2\,\sigma t_{0;23}$. As $\sigma$ runs
over $S_4$, these run over all incident and all non-incident basis
vectors up to sign, so the $24$ vectors $\sigma\xi$ span $T$, which
has dimension $24$; they are a basis. $\Phi$ sends the basis
$\{\sigma\}$ of $\mathbb{R}[S_4]$ to this basis, and
$\Phi(\rho\sigma) = \rho\sigma\xi = \rho\,\Phi(\sigma)$.
\hfill $\square$

\paragraph{Corollary 1.} \textit{Each irreducible representation
occurs in $T$ with multiplicity equal to its dimension,}
\begin{equation}
  T \cong \mathbf{1} \oplus \sgn \oplus 2\,[2,2] \oplus 3\,\std
  \oplus 3\,(\sgn \otimes \std),
\end{equation}
\textit{and $T = T_{\mathrm{inc}} \oplus T_{\mathrm{non}}$ with
$T_{\mathrm{inc}} \cong \Ind_H^{S_4}\mathbf{1}$ and
$T_{\mathrm{non}} \cong \Ind_H^{S_4}\varepsilon$, where
$\varepsilon$ is the sign of $H$.}

\paragraph{Proof.} The first statement holds for every regular
representation~\cite{Serre}. A module spanned by one orbit of lines,
on which the stabiliser $H$ of one line acts by a character $\chi$,
is $\Ind_H^{S_4}\chi$~\cite{Serre}; apply the Lemma. As a check,
$\Ind_H(\mathbf{1} \oplus \varepsilon) = \Ind_H \mathbb{R}[H] =
\mathbb{R}[S_4]$ by transitivity of induction, and counting fixed
lines with their signs gives the character of $T$ as
$(24, 0, 0, 0, 0)$ on the classes $e$, $(ab)$, $(ab)(cd)$, $(abc)$,
$(abcd)$. \hfill $\square$

% ----------------------------------------------------------------------

\section*{The Sign Representation}

\paragraph{Theorem 2.} \textit{$\sgn$ occurs exactly once in $T$
and not at all in $C$. In $T$ it is spanned by $\omega$, with
components $\omega^a{}_{bc} = \sum_d \varepsilon_{abcd}$, which lies
in $T_{\mathrm{non}}$.}

\paragraph{Proof.} By Frobenius reciprocity~\cite{Serre}, $\sgn$
occurs in $\Ind_H\chi$ once if $\chi = \sgn|_H$ and otherwise not.
Since $\sgn|_H = \varepsilon$, it occurs once in
$T_{\mathrm{non}}$ and not in $T_{\mathrm{inc}}$. The basis
$f_a \otimes f_b f_c$ of $C$ is permuted without signs, in four
orbits: $a = b = c$ ($4$), $b = c \neq a$ ($12$),
$a \in \{b,c\}$ with $b \neq c$ ($12$), and $a, b, c$ distinct
($12$). Each stabiliser contains a transposition acting by $+1$, so
$C$ is a sum of modules $\Ind_{H_i}\mathbf{1}$ with
$\sgn|_{H_i} \neq \mathbf{1}$, none of which contains $\sgn$.
Finally $\varepsilon_{\sigma(a)\sigma(b)\sigma(c)\sigma(d)} =
\sgn(\sigma)\,\varepsilon_{abcd}$ and the sum over $d$ is
invariant, so $\sigma\omega = \sgn(\sigma)\,\omega$; $\omega \neq 0$
is antisymmetric in $b, c$ and vanishes unless $a, b, c$ are
distinct. \hfill $\square$

% ----------------------------------------------------------------------

\section*{Vector, Axial and Tensor Parts}

Lower indices with $|\eta| = I$. The torsion splits into its vector
part $T_v = \{\delta^a{}_b v_c - \delta^a{}_c v_b\}$, its axial part
$T_{ax} = \{\varepsilon_{abcd}\, w^d\}$, both of dimension $4$, and
its tensor part $T_t$, of dimension $16$, the torsion that is
trace-free and has no totally antisymmetric part~\cite{Hehl1976}.

\paragraph{Theorem 3.} \textit{$T_v \subset T_{\mathrm{inc}}$ and
$T_{ax} \subset T_{\mathrm{non}}$, and
$T_t = (T_{\mathrm{inc}} \ominus T_v) \oplus
(T_{\mathrm{non}} \ominus T_{ax})$. As $S_4$-modules}
\begin{equation}
\begin{aligned}
  T_v &\cong V = \mathbf{1} \oplus \std, \\
  T_{ax} &\cong \sgn \otimes V = \sgn \oplus (\sgn \otimes \std),
\end{aligned}
\end{equation}
\textit{and each half of $T_t$ is
$[2,2] \oplus \std \oplus (\sgn \otimes \std)$. Lowered by $\eta$
itself, the axial part is not $S_4$-invariant.}

\paragraph{Proof.} The components of $T_v$ vanish unless
$a \in \{b,c\}$, and those of $T_{ax}$ unless $a, b, c$ are
distinct. The trace $T^a{}_{ac}$ involves only incident components,
and $T_v$ is the image of its adjoint, so $T_{\mathrm{inc}} \ominus
T_v$ is the trace-free part of $T_{\mathrm{inc}}$ and has no totally
antisymmetric part. Likewise $T_{ax}$ is the image of total
antisymmetrisation, which involves only non-incident components, so
$T_{\mathrm{non}} \ominus T_{ax}$ is trace-free with no totally
antisymmetric part. Both halves lie in $T_t$, and their dimensions
sum to $16 = \dim T_t$. The maps $v \mapsto T_v$ and
$w \mapsto T_{ax}$ are equivariant, the second with the sign of
$\varepsilon$. Counting fixed lines with signs gives
$\chi_{\mathrm{inc}} = (12, 2, 0, 0, 0)$ and
$\chi_{\mathrm{non}} = (12, -2, 0, 0, 0)$; subtracting
$\chi_V = (4, 2, 0, 1, 0)$ and $\sgn\,\chi_V$ leaves
$(8, 0, 0, -1, 0)$ for each half, the character of
$[2,2] \oplus \std \oplus (\sgn \otimes \std)$. With $\eta$, the
axial tensor whose $\eta$-lowered components are
$\varepsilon_{abc2}$ is sent by $(01)$ to one with lowered components
$T_{013} = -1$ and $T_{301} = +1$, which is not totally
antisymmetric. \hfill $\square$

% ----------------------------------------------------------------------

\section*{Why Four Legs}

\paragraph{Proposition.} \textit{For $K$ legs,
$\dim T = K\binom{K}{2}$ equals $|S_K| = K!$ only for $K = 2$ and
$K = 4$, and at both the torsion sector is the regular
representation.}

\paragraph{Proof.} For $K \ge 2$ the equation reduces to
$K/2 = (K-2)!$. It holds at $K = 2$ and $K = 4$, fails at $K = 3$,
and for $K \ge 5$,
$(K-2)! \ge (K-2)(K-3) \ge 2(K-2) > K/2$. At $K = 1$,
$\dim T = 0 \neq 1$. At $K = 2$, $T$ is spanned by $t_{0;01}$ and
$t_{1;01}$, which the transposition exchanges with a sign, so its
character is $(2, 0)$; $K = 4$ is Theorem 1. \hfill $\square$

% ----------------------------------------------------------------------

\section*{Scope}

$\sgn(\sigma) = \det P_\sigma$: the sign representation is the
orientation character of $GL(V)$ restricted to permutations of the
legs. Spatial parity, $\mathrm{diag}(1,-1,-1,-1)$, also reverses
orientation, but it is not a permutation of the legs, and these
theorems do not describe its action. Theorems 1 and 2 hold for any
four-index frame, with or without a metric. Theorem 3 uses the
magnitude reading $|\eta| = I$; with the sign of $\eta$ kept, the
axial part is not $S_4$-invariant.

% ----------------------------------------------------------------------

\section*{Summary}

The $24$ torsion components of the tetrad connection are one copy
of the group algebra of the legs' permutations, generated by one
explicit tensor. The sign representation lives in the torsion alone,
as its totally antisymmetric part along $(1,1,1,1)$. Read by
magnitude, the vector part of the torsion is incident, the axial
part is non-incident, and the tensor part is two isomorphic halves.
General relativity's torsion-free condition therefore removes a
complete copy of $\mathbb{R}[S_4]$, and with it the only scalar
response of the connection to odd permutations of the legs.

% ----------------------------------------------------------------------

\enlargethispage{2\baselineskip}
\begin{thebibliography}{9}

\bibitem{Nelson2026frame}
S.~Nelson,
\textit{The Spatial Geometry of the Lorentz Tetrad} (2026).

\bibitem{Hehl1976}
F.~W.~Hehl, P.~von der Heyde, G.~D.~Kerlick and J.~M.~Nester,
\textit{General relativity with spin and torsion: foundations and
prospects},
Rev.\ Mod.\ Phys.\ \textbf{48} (1976) 393.

\bibitem{Serre}
J.-P.~Serre,
\textit{Linear Representations of Finite Groups}
(Springer, New York, 1977).

\bibitem{FultonHarris}
W.~Fulton and J.~Harris,
\textit{Representation Theory: A First Course}
(Springer, New York, 1991).

\end{thebibliography}

\end{document}
